\section{Complete definition and entry-point index}\label{app:correspondence}
The two default entry modules in \texttt{lakefile.toml} are
\path{CDCert.Main} and \path{CDCert.TwoInputs}. The first exports seven
results and the second three. The following list includes every one
of these ten declarations separately, even when a wrapper repeats a
supporting result. It then lists all ten additional results in the
independent axiom audit. Throughout the table, \(n=k+1\) for domain
statements on \(\mathrm{Fin}(k+1)\).

The companion \path{audit/TUTORIAL_LEAN_MAP.json} gives the exact
source file, line, declaration signature, tutorial location and
normalisation for every entry. \path{audit/ENTRYPOINT_SIGNATURES.txt}
contains the literal source signatures; the earlier compiler-printed
statements and axiom lists are preserved separately. The names below
are declarations, not a list inferred only from file comments.

\subsection{All ten final exported declarations}
\begin{longtable}{@{}>{\raggedright\arraybackslash}p{0.43\linewidth}>{\raggedright\arraybackslash}p{0.52\linewidth}@{}}
\toprule Entry and Lean declaration & Exact mathematical content and hypothesis\\\midrule
\endhead
L01 \path{CDCert.Main.general_recurrence} & $n\geq3$: $(11n+8)f(n)\leq32nf(n-1)$. See Theorem~\ref{thm:final}.\\[0.65em]
L02 \path{CDCert.Main.general_geometric} & $n\geq3$: $f(n)\leq4(32/11)^{n-3}$. See Theorem~\ref{thm:final}.\\[0.65em]
L03 \path{CDCert.Main.peak_recurrence} & $n\geq3$: $3h(n)\leq8h(n-1)$. See Theorem~\ref{thm:final}.\\[0.65em]
L04 \path{CDCert.Main.peak_geometric} & $n\geq3$: $h(n)\leq4(8/3)^{n-3}$. See Theorem~\ref{thm:final}.\\[0.65em]
L05 \path{CDCert.Main.general_limit} & No extra premises: $\exists L_f$, $f(n)^{1/n}\to L_f\leq32/11$. See Theorem~\ref{thm:final}.\\[0.65em]
L06 \path{CDCert.Main.peak_limit} & No extra premises: $\exists L_h$, $h(n)^{1/n}\to L_h\leq8/3$. See Theorem~\ref{thm:final}.\\[0.65em]
L07 \path{CDCert.Main.general_polynomial_correction} & No extra premises: $f(n)=O(n^{-8/11}(32/11)^n)$. See Theorem~\ref{thm:final}.\\[0.65em]
L08 \path{CDCert.TwoInputs.first_input} & $D$ maximal, $n=k+1$: $G+(n+1)M\leq3\Sigma$. See Theorem~\ref{thm:gap}.\\[0.65em]
L09 \path{CDCert.TwoInputs.second_input} & $D$ full, $n=k+1$: $3nM\leq8(\Sigma+G)$. See Theorem~\ref{thm:second}.\\[0.65em]
L10 \path{CDCert.TwoInputs.transfer} & $D$ maximal, $n=k+1$: $(11n+8)M\leq32\Sigma$. See Theorem~\ref{thm:transfer}.\\[0.65em]
\bottomrule\end{longtable}

\subsection{The ten additional audited supporting results}
\begin{longtable}{@{}>{\raggedright\arraybackslash}p{0.43\linewidth}>{\raggedright\arraybackslash}p{0.52\linewidth}@{}}
\toprule Entry and Lean declaration & Exact mathematical content and hypothesis\\\midrule
\endhead
L11 \path{CDCert.Recurrences.gap_inequality} & $D$ maximal: $G+(n+1)M\leq3\Sigma$. See Theorem~\ref{thm:gap}.\\[0.65em]
L12 \path{CDCert.Recurrences.deletion_lower_bound} & $D$ maximal: $(11n+8)M\leq32\Sigma$. See Theorem~\ref{thm:transfer}.\\[0.65em]
L13 \path{CDCert.Bridge.InF.density} & $K\in\mathcal F_n$, $n=k+1$: $3n|K|\leq8\Sigma(K)$. See Theorem~\ref{thm:density}.\\[0.65em]
L14 \path{CDCert.FreeLocked.free_le_locked} & $D$ maximal: $\mathrm{fr}(D)\leq\mathrm{lk}(D)$. See Proposition~\ref{prop:injection}.\\[0.65em]
L15 \path{CDCert.K6.finite_lemma} & $\texttt{DomE}(7,J)$: the integer-scaled seven-window inequality. See Theorem~\ref{thm:finite}.\\[0.65em]
L16 \path{CDCert.K6.adj_le} & $\texttt{IsDom}(n,K)$: $8A(K)\leq5n|K|$. See Theorem~\ref{thm:density}.\\[0.65em]
L17 \path{CDCert.Supermult.f_supermul} & All $m,n\in\mathbb N$: $f(m)f(n)\leq f(m+n)$. See Theorem~\ref{thm:final}.\\[0.65em]
L18 \path{CDCert.Supermult.h_supermul} & All $m,n\in\mathbb N$: $h(m)h(n)\leq h(m+n)$. See Theorem~\ref{thm:final}.\\[0.65em]
L19 \path{CDCert.Recurrences.general_recurrence} & All $k\in\mathbb N$: $(11(k+1)+8)f(k+1)\leq32(k+1)f(k)$. See Theorem~\ref{thm:final}.\\[0.65em]
L20 \path{CDCert.Bounds.peakPit_recurrence} & All $k\in\mathbb N$: $3h(k+1)\leq8h(k)$. See Theorem~\ref{thm:final}.\\[0.65em]
\bottomrule\end{longtable}

\subsection{Definitions and representation conventions}
The following groups identify the mathematical objects whose meanings
must match. Their full source locations and signatures, together with
all named helper lemmas beside the proofs, are in the machine-readable map.
\begin{longtable}{@{}>{\raggedright\arraybackslash}p{0.32\linewidth}>{\raggedright\arraybackslash}p{0.63\linewidth}@{}}
\toprule Tutorial object & Lean definitions and precise convention\\\midrule\endhead
Orders and domains & \path{Condorcet.IsStrictTotal}, \path{Order}, \path{Domain}: Boolean strict total relations, finite sets of distinct orders.\\[0.5em]
Majority and Condorcet & \path{Condorcet.majority}, \path{IsCondorcet}: strict majority in every odd profile; repetitions permitted.\\[0.5em]
Never-laws & \path{Condorcet.Never}, \path{ValueRestricted}, \path{ExtremeRestricted}: forbidden positions 0,1,2; extremes 0 and 2.\\[0.5em]
Peak-pit and extrema & \path{Condorcet.IsPeakPit}, \path{f}, \path{h}: semantic domains and actual finite maxima, with no shared axis.\\[0.5em]
Full and maximal & \path{Condorcet.IsFull}, \path{IsMaximal}: full solutions of one selected law per triple, versus inclusion-maximal semantic domain.\\[0.5em]
Extreme full class & \path{Condorcet.IsExtremeFull}, \path{InF}: extreme laws plus optional fixed pair comparisons; only \texttt{InF} requires nonemptiness.\\[0.5em]
Deletions and gaps & \path{Condorcet.deleteDomain}, \path{deletionSum}, \path{insertionWord}, \path{gapCount}; one separation per pair of consecutive runs, not one per absent slot.\\[0.5em]
Adjacency and moves & \path{Condorcet.adjacencyCount}, \path{HasLift}, \path{HasDrop}, \path{Free}, \path{Frozen}, \path{Locked}: each swap edge twice; lifts/drops may jump.\\[0.5em]
List domains & \path{CDCert.K6.IsDom}, \path{DomE}, \path{CDCert.ListOrder.groupsEP}: full constraint systems on all permutations, including pair choices; values on other lists are irrelevant.\\[0.5em]
Finite potentials & \path{CDCert.K6.cost10}, \path{qD}, \path{lamD}: respectively \(10C,2dq,2d\Lambda\), with fixed literal data.\\
\bottomrule\end{longtable}

\subsection{What is not an additional formal theorem here}
The helper facts about a particular labelled example, the arbitrary
number-of-copies deletion-product formula, the informal twelve-idea
classification and the discovery narrative are written explanations.
They are not advertised as separately exported Lean theorems. The
formal amplification instead uses the proved adjacency product and
powers of two. The original project's arbitrary-real-density transfer
has no separate parameterised declaration; the proved \(3/8\) instance
is exactly the one this tutorial uses.

There are 326 project source modules, with 319 in the dependency
closure of the two final entry modules. Seven auxiliary modules lie
outside that closure: \path{Audit}, \path{CDCert},
\path{CDCert.Levels}, \path{ClaudeAudit}, \path{Condorcet},
\path{KernelTransfer}, \path{ThreeOrders}. They are supplied and were
compiled in the earlier complete audit. They are not extra assumptions
or missing entry points of the final theorem. Generated finite-family
theorems and proof-engineering helpers are included as dependencies;
their precise statements can be inspected in the complete source.
